% Options for packages loaded elsewhere
\PassOptionsToPackage{unicode}{hyperref}
\PassOptionsToPackage{hyphens}{url}
\PassOptionsToPackage{dvipsnames,svgnames,x11names}{xcolor}
\documentclass[
  10pt,
  fontset=fandol]{ctexart}
\usepackage{xcolor}
\usepackage[margin=16mm]{geometry}
\usepackage{amsmath,amssymb}
\setcounter{secnumdepth}{-\maxdimen} % remove section numbering
\usepackage{iftex}
\ifPDFTeX
  \usepackage[T1]{fontenc}
  \usepackage[utf8]{inputenc}
  \usepackage{textcomp} % provide euro and other symbols
\else % if luatex or xetex
  \usepackage{unicode-math} % this also loads fontspec
  \defaultfontfeatures{Scale=MatchLowercase}
  \defaultfontfeatures[\rmfamily]{Ligatures=TeX,Scale=1}
\fi
\usepackage{lmodern}
\ifPDFTeX\else
  % xetex/luatex font selection
\fi
% Use upquote if available, for straight quotes in verbatim environments
\IfFileExists{upquote.sty}{\usepackage{upquote}}{}
\IfFileExists{microtype.sty}{% use microtype if available
  \usepackage[]{microtype}
  \UseMicrotypeSet[protrusion]{basicmath} % disable protrusion for tt fonts
}{}
\makeatletter
\@ifundefined{KOMAClassName}{% if non-KOMA class
  \IfFileExists{parskip.sty}{%
    \usepackage{parskip}
  }{% else
    \setlength{\parindent}{0pt}
    \setlength{\parskip}{6pt plus 2pt minus 1pt}}
}{% if KOMA class
  \KOMAoptions{parskip=half}}
\makeatother
\setlength{\emergencystretch}{3em} % prevent overfull lines
\providecommand{\tightlist}{%
  \setlength{\itemsep}{0pt}\setlength{\parskip}{0pt}}
\usepackage{amsmath,amssymb,mathtools,needspace}
\xeCJKDeclareCharClass{CJK}{"2032,"0400->"04FF,"2160->"216B,"2460->"2473,"25A0,"25C7,"25CB,"25CF}
\IfFontExistsTF{Arial Unicode MS}{\xeCJKsetup{AutoFallBack=true}\setCJKfallbackfamilyfont{\CJKrmdefault}{Arial Unicode MS}}{}
\setcounter{secnumdepth}{0}
\setlength{\emergencystretch}{2em}
\allowdisplaybreaks
\usepackage{bookmark}
\IfFileExists{xurl.sty}{\usepackage{xurl}}{} % add URL line breaks if available
\urlstyle{same}
\hypersetup{
  pdftitle={只要做一次``俯冲''},
  colorlinks=true,
  linkcolor={Maroon},
  filecolor={Maroon},
  citecolor={Blue},
  urlcolor={Blue},
  pdfcreator={LaTeX via pandoc}}

\title{只要做一次``俯冲''}
\author{}
\date{}

\begin{document}
\maketitle

\subsubsection{12.3.3　只要做一次``俯冲''}\label{ux53eaux8981ux505aux4e00ux6b21ux4fefux51b2}

偏差比几乎是个定值，基于这个奇妙的规律，分析误差只是一蹴而就的事，据此只要做一次``俯冲''便能捕捉到所要的校正公式。

事实上，由于在割圆计算过程中偏差比近似等于 4，从而得出一系列近似关系

\href{../../source-images/pdf-300.jpeg}{原页图像}

式：

\[
\begin{aligned}
S_{2n}-S_n&\approx4(S_{4n}-S_{2n}),\\
S_{4n}-S_{2n}&\approx4(S_{8n}-S_{4n}),\\
S_{8n}-S_{4n}&\approx4(S_{16n}-S_{8n}),\\
&\vdots\\
S_{2N}-S_N&\approx4(S_{4N}-S_{2N}),
\end{aligned}
\]

式中，\(N\) 是某个远大于 \(n\) 的正整数。

将这些式子累加在一起，其中间项相互抵消，得

\[
S_{2N}-S_n\approx4(S_{4N}-S_{2n}).
\]

这样，若取 \(S_{4N}\) 和 \(S_{2N}\) 作为 \(S_{2n}\) 的校正值 \(\widehat S\)，则有

\[
\widehat S-S_n=4(\widehat S-S_{2n}),
\tag{12.3.2}
\]

即有

\[
\frac{\widehat S-S_n}{\widehat S-S_{2n}}\approx4.
\]

如果称近似值 \(S_n\) 与校正值 \(\widehat S\) 两者之差为\textbf{残差}，则上式说明，\textbf{若偏差比几乎为定值 4，那么残差比也近似等于 4。}

这样一来，精加工方法的设计就水到渠成了。事实上，从式（12.3.2）解出未知的 \(\widehat S\)，有

\[
\widehat S=S_{2n}+\frac{1}{3}(S_{2n}-S_n).
\tag{12.3.3}
\]

这个式子的含义是，在偏差比``几乎''为定值 4 的情况下，近似值的残差比也近似等于 4，因而残差 \(\widehat S-S_{2n}\) 几乎等于偏差 \(S_{2n}-S_n\) 的 \(\dfrac13\)，从而立即导出所求的加速公式（12.3.3）。

值得指出的是，刘徽神算（12.2.2）即

\[
\widehat S=S_{192}+\frac{36}{105}(S_{192}-S_{96})
\]

中的校正因子 \(\omega=\dfrac{36}{105}\) 非常接近实际精加工公式（12.3.3）中的校正因子 \(\omega=\dfrac13=\dfrac{35}{105}\)。至于刘徽为什么这样处理，\textbf{一个明显的理由是，刘徽过于偏爱计算公式与计算结果的简洁美。}

\end{document}
